\documentclass[11pt,a4paper]{article}

\usepackage[T1]{fontenc}
\usepackage[utf8]{inputenc}
\usepackage{geometry}
\usepackage{xcolor}
\usepackage{hyperref}
\usepackage{longtable}
\usepackage{booktabs}
\usepackage{listings}
\usepackage{fancyhdr}

\geometry{margin=1in}
\hypersetup{
  colorlinks=true,
  linkcolor=blue,
  urlcolor=blue,
  pdftitle={Raden Language Manual},
  pdfauthor={Raden project}
}

\definecolor{codebg}{RGB}{248,248,248}
\definecolor{codeframe}{RGB}{220,220,220}

\lstdefinelanguage{Raden}{
  morekeywords={
    let,set,const,unlet,defun,apply,demac,call,pcall,ret,if,else,end,loop,break,
    continue,load,loadnative,add_load_path,add_native_path,true,false,null,
    print,type,to_string,exit,append,remove,index,len,enum,reset,error,assert,
    check,scall,module,open,match
  },
  sensitive=true,
  morecomment=[s]{[*}{*]},
  morestring=[b]"
}

\lstset{
  basicstyle=\ttfamily\small,
  backgroundcolor=\color{codebg},
  frame=single,
  rulecolor=\color{codeframe},
  columns=fullflexible,
  keepspaces=true,
  breaklines=true,
  showstringspaces=false,
  tabsize=2
}

\setlength{\headheight}{14pt}
\pagestyle{fancy}
\fancyhf{}
\lhead{Raden}
\rhead{Language Manual}
\cfoot{\thepage}

\newcommand{\raden}{\texttt{raden}}
\newcommand{\rdn}{\texttt{rdn}}
\newcommand{\code}[1]{\texttt{#1}}

\title{Raden Language Manual}
\author{Raden project}
\date{}

\begin{document}
\maketitle

\begin{center}
Raden is a small stack-based scripting language with a C interpreter,
standard libraries, and a native module API.
\end{center}

\tableofcontents
\newpage

\section{Introduction}

Raden is built around a value stack. Values are pushed first, then
operators and builtins consume values from the stack and push results
back. This makes the language postfix:

\begin{lstlisting}[language=Raden]
2 3 + print
\end{lstlisting}

This program pushes \code{2}, pushes \code{3}, applies \code{+}, and
prints the result. Operand order matters. The expression \code{i 10 <}
means \code{i < 10}; the expression \code{10 i <} means \code{10 < i}.

The interpreter is written in C. The public embedding entry point is:

\begin{lstlisting}[language=C]
int rdn_main(int argc, char **argv);
\end{lstlisting}

\section{Quick Start}

\subsection{Install}

\begin{lstlisting}[language=bash]
curl -fsSL https://raw.githubusercontent.com/abdorayden/rdn/master/install | sh
\end{lstlisting}

Or from a cloned checkout:

\begin{lstlisting}[language=bash]
./install
\end{lstlisting}

By default, libraries install to \code{/usr/local/share/rdn/libs} and
native modules install to \code{/usr/local/share/rdn/nativelibs}.

\subsection{Build}

\begin{lstlisting}[language=bash]
git clone https://github.com/abdorayden/rdn.git
cd rdn
make
./main script.rdn
\end{lstlisting}

\subsection{REPL}

Run \code{./main} with no script path to start the REPL:

\begin{lstlisting}[language=bash]
$ ./main
raden repl
press Ctrl-D to exit
RDN >> 2 3 + print
\end{lstlisting}

\subsection{Tests}

\begin{lstlisting}[language=bash]
bash test/run.sh
\end{lstlisting}

\section{Values and Syntax}

\subsection{Value Types}

\begin{longtable}{lll}
\toprule
Type & Syntax & Example \\
\midrule
Integer & \code{123} & \code{42} \\
Double & \code{123.456} & \code{3.14} \\
String & \code{"..."} & \code{"hello"} \\
Boolean & \code{true}, \code{false} & \code{true} \\
Null & \code{null} & \code{null} \\
List & \code{(...)} & \code{(1 2 "three")} \\
Identifier & bare word & \code{x}, \code{hello} \\
\bottomrule
\end{longtable}

\subsection{Numbers}

\begin{lstlisting}[language=Raden]
10
2.5
\end{lstlisting}

\subsection{Strings}

Strings use double quotes. Common escapes include \code{\textbackslash n},
\code{\textbackslash t}, \code{\textbackslash "}, \code{\textbackslash\textbackslash},
and \code{\textbackslash e}.

\begin{lstlisting}[language=Raden]
"hello"
"line1\nline2"
"quote: \""
\end{lstlisting}

\subsection{Booleans and Null}

\begin{lstlisting}[language=Raden]
true
false
null
\end{lstlisting}

\subsection{Lists}

Lists use parentheses:

\begin{lstlisting}[language=Raden]
(1 2 3)
("hello" 4 true)
\end{lstlisting}

List literals execute their contents. Values left on the temporary list
stack become list items. This means conditionals, loops, builtins, and
function calls can run inside a list literal:

\begin{lstlisting}[language=Raden]
(2 3 1 true if 10 end)

0 i let
(i 10 < loop
  i
  i 1 + i set
  i 10 <
end) print
\end{lstlisting}

\subsection{Comments}

\begin{lstlisting}[language=Raden]
[* this is a comment *]
42 [* inline comment *] print
\end{lstlisting}

\section{Operators}

Operators consume operands from the stack.

\begin{longtable}{ll}
\toprule
Operator & Meaning \\
\midrule
\code{+} & addition \\
\code{-} & subtraction \\
\code{*} & multiplication \\
\code{/} & division \\
\code{=} & equality \\
\code{!=} & inequality \\
\code{<} & less than \\
\code{>} & greater than \\
\code{<=} & less than or equal \\
\code{>=} & greater than or equal \\
\code{!} & boolean not \\
\code{\&} & boolean and, or integer bitwise and \\
\code{|} & boolean or, or integer bitwise or \\
\code{\^{}} & integer bitwise xor \\
\code{<<} & integer left shift \\
\code{>>} & integer right shift \\
\bottomrule
\end{longtable}

\begin{lstlisting}[language=Raden]
2 3 + print
10 3 - print
5 3 > print
true false & print
4 1 << print
\end{lstlisting}

\section{Variables}

\subsection{let}

\code{let} binds a value to a mutable name:

\begin{lstlisting}[language=Raden]
10 x let
x print
\end{lstlisting}

\subsection{set}

\code{set} updates an existing binding:

\begin{lstlisting}[language=Raden]
x 1 + x set
x print
\end{lstlisting}

\subsection{const}

\code{const} creates an immutable binding. Reassigning a constant is an
error:

\begin{lstlisting}[language=Raden]
100 limit const
limit print
\end{lstlisting}

\subsection{unlet}

\code{unlet} removes a binding:

\begin{lstlisting}[language=Raden]
10 x let
x unlet
\end{lstlisting}

\subsection{enum and reset}

\begin{lstlisting}[language=Raden]
enum ONE const
enum TWO const
enum THREE const
reset
\end{lstlisting}

\subsection{Built-in Environment Variables}

\begin{longtable}{ll}
\toprule
Variable & Value \\
\midrule
\code{\_\_host\_os} & host OS name, such as \code{"linux"} \\
\code{\_\_sharedlib\_ext} & native library extension, such as \code{".so"} \\
\code{\_\_argv} & script path followed by extra CLI arguments \\
\code{\_\_line\_col} & source location helper for the called variable \\
\bottomrule
\end{longtable}

\subsection{Identifier Semantics}

Identifiers are context-sensitive. Names used before \code{let},
\code{set}, \code{const}, \code{defun}, or \code{call} are treated as name
operands. Normal variable lookup may clone the underlying value, while
list and string variables are often preserved as named targets for
mutation-oriented builtins such as \code{append}, \code{remove},
\code{index}, and \code{len}.

When a script needs a copied list value before mutation, \code{dup} may be
required:

\begin{lstlisting}[language=Raden]
next dup current set pop
\end{lstlisting}

\section{Functions}

\subsection{defun}

\begin{lstlisting}[language=Raden]
double defun
  2 *
end

5 double call print
\end{lstlisting}

\subsection{apply}

\code{apply} defines a function-like word that is invoked directly,
without \code{call}:

\begin{lstlisting}[language=Raden]
triple apply
  3 *
end

4 triple print
\end{lstlisting}

\subsection{demac}

\code{demac} defines a textual macro. When the macro word is read, its
body is expanded into the active source stream before parsing continues.
This allows aliases for syntax words, not only value-producing code.

\begin{lstlisting}[language=Raden]
for demac loop end

0 i let
true for
  i print "\n" print
  i 1 + i set
  i 3 <
end
\end{lstlisting}

\subsection{call}

\code{call} invokes a function by name. A string can also name the
function:

\begin{lstlisting}[language=Raden]
"double" call call
\end{lstlisting}

\subsection{pcall}

\code{pcall} performs a protected call and catches runtime errors. On
success the stack receives \code{result true}; on error it receives
\code{"error message" false}.

\begin{lstlisting}[language=Raden]
risky defun
  null 0 index
end

risky pcall
if
  "success: " print print "\n" print
else
  "error: " print print "\n" print
end
\end{lstlisting}

\subsection{ret}

\code{ret} stops execution of the current function immediately. Raden
uses stack returns, so any values already left on the stack are the
function's return values.

\begin{lstlisting}[language=Raden]
check-less-three defun
  n let
  n 3 < if
    true
    ret
  end
  false
end
\end{lstlisting}

\subsection{Scope}

Functions create a new variable scope. Bindings created inside a function
are local unless \code{set} updates an existing outer binding.

\section{Modules}

\code{module} groups functions, macros, apply-words, and variables into a
named namespace. Names inside a module are prefixed with \code{ModuleName::}
for functions and \code{ModuleName.} for variables to distinguish them from
top-level names.

\subsection{module}
\begin{lstlisting}[language=Raden]
MyModule module
    func1 defun
        "hello\n" print
    end
    app1 apply
        "world\n" print
    end
end

MyModule::func1 call
MyModule::app1
\end{lstlisting}

Variables and constants inside a module also use \code{::} as prefix:
\begin{lstlisting}[language=Raden]
Cfg module
    10 limit let
end

Cfg::limit print
\end{lstlisting}

\subsection{open}

\code{open} exposes all members of a module in the current scope without
requiring the module prefix. Unqualified names are created for each
function and variable.

\begin{lstlisting}[language=Raden]
MyModule open

[* now call without prefix *]
func1 call
app1
\end{lstlisting}

For functions, \code{open} creates short aliases using \code{::}
prefix matching; for variables, it uses \code{.} prefix matching.

\subsection{Nested Modules}

Modules can be nested. Inner module names are chained with \code{::}:

\begin{lstlisting}[language=Raden]
Mod1 module
    NestedMod1 module
        nm apply
            "nm\n" print
        end
    end
end

Mod1::NestedMod1::nm
\end{lstlisting}

\subsection{Cross-file Modules}

Modules work across files. A module defined in one file can be accessed
from another file by loading it first:

\begin{lstlisting}[language=Raden]
[* file: module.rdn *]
MyModule module
    greet defun
        "hi\n" print
    end
end

[* file: main.rdn *]
"./module.rdn" load
MyModule::greet call
MyModule open
greet call
\end{lstlisting}

\section{Control Flow}

\subsection{if}

\code{if} consumes one boolean. \code{else} is optional. Blocks terminate
with \code{end}.

\begin{lstlisting}[language=Raden]
true if
  "yes" print
else
  "no" print
end
\end{lstlisting}

\subsection{loop}

\code{loop} consumes a boolean condition at entry. Each iteration must
leave the next boolean condition on the stack.

\begin{lstlisting}[language=Raden]
0 i let
i 5 < cond let

cond loop
  i print "\n" print
  i 1 + i set
  i 5 < cond set
  cond
end
\end{lstlisting}

\subsection{break, continue, and exit}

\begin{lstlisting}[language=Raden]
0 i let
true loop
  i 10 > if break end
  i 5 = if i 1 + i set continue end
  i print "\n" print
  i 1 + i set
  true
end

0 exit
\end{lstlisting}

\section{Builtins}

\begin{longtable}{lll}
\toprule
Builtin & Stack shape & Description \\
\midrule
\code{print} & \code{value ->} & Print to stdout \\
\code{pop} & \code{value ->} & Discard top value \\
\code{dup} & \code{value -> value value} & Duplicate top value \\
\code{swap} & \code{a b -> b a} & Swap top two values \\
\code{type} & \code{value -> string} & Return type name \\
\code{to\_string} & \code{value -> string} & Convert to string \\
\code{assert} & \code{condition message ->} & Raise message if condition is false \\
\code{let} & \code{value name ->} & Create mutable binding \\
\code{set} & \code{value name ->} & Update binding \\
\code{const} & \code{value name ->} & Create constant binding \\
\code{unlet} & \code{name ->} & Remove binding \\
\code{enum} & \code{-> integer} & Push next enum integer \\
\code{reset} & \code{->} & Reset enum counter \\
\code{index} & \code{target index -> value} & Index list or string \\
\code{append} & \code{target value -> target} & Append to list or string \\
\code{remove} & \code{target index -> target} & Remove item or character \\
\code{len} & \code{target -> integer} & Length of list or string \\
\code{load} & \code{path ->} & Execute a Raden source file \\
\code{loadnative} & \code{path ->} & Load a native shared library \\
\code{add\_load\_path} & \code{path ->} & Add script search path \\
\code{add\_native\_path} & \code{path ->} & Add native search path \\
\code{defun} & \code{name ... end ->} & Define named function \\
\code{apply} & \code{name ... end ->} & Define direct-call word \\
\code{demac} & \code{name tokens... end ->} & Define read-time macro \\
\code{call} & \code{name args... -> results...} & Invoke a function \\
\code{pcall} & \code{name args... -> result true | error false} & Protected call \\
\code{ret} & \code{values... -> values...} & Return from current function \\
\code{error} & \code{string -> error} & Raise an interpreter error \\
\code{if} & \code{cond ... end ->} & Conditional execution \\
\code{loop} & \code{cond ... next-cond end ->} & Loop while true \\
\code{break} & \code{->} & Exit loop \\
\code{continue} & \code{->} & Continue loop \\
\code{exit} & \code{status ->} & Exit process \\
\bottomrule
\end{longtable}

\subsection{String Matching}

\code{match} performs glob-style pattern matching and returns a boolean:

\begin{lstlisting}[language=Raden]
"hello.txt" "*.txt" match print   [* true *]
"hello.md" "*.txt" match print    [* false *]
"abc" "a?c" match print           [* true *]
\end{lstlisting}

Pattern syntax:
\begin{itemize}
\item \code{*} — matches any sequence of characters (including empty)
\item \code{?} — matches any single character
\item \code{[abc]} — matches any character in the set
\item \code{[!abc]} or \code{[\^{}abc]} — matches any character not in the set
\item \code{[a-z]} — matches any character in the range
\end{itemize}

\section{Lists and Strings}

\begin{lstlisting}[language=Raden]
(10 20 30) 1 index print
(1 2) 3 append print
(1 2 3) 1 remove print
(1 2 3) len print
\end{lstlisting}

The same sequence builtins work on strings:

\begin{lstlisting}[language=Raden]
"hello" 1 index print
"hello" 1 remove print
"hi" len print

"hi" text let
text "!" append
text print
\end{lstlisting}

\section{Error Handling}

\code{pcall} is the primary error-handling mechanism. It invokes a
function and catches errors raised during execution.

\begin{longtable}{ll}
\toprule
Situation & Stack after \code{pcall} \\
\midrule
Success & \code{... function\_result true} \\
Error caught & \code{... "error message" false} \\
\bottomrule
\end{longtable}

Errors that happen before the function executes, such as a missing
function name, propagate normally.

\section{Loading Code}

\subsection{Raden Scripts}

\begin{lstlisting}[language=Raden]
"./libs/math.rdn" load
4 sqrt call print
\end{lstlisting}

Paths are resolved relative to the currently running source file and then
through configured search paths.

\subsection{Native Libraries}

\begin{lstlisting}[language=Raden]
"../nativelibs/files.so" loadnative
"./Makefile" readLines call print
\end{lstlisting}

Use \code{\_\_sharedlib\_ext} for cross-platform native library paths.

\begin{lstlisting}[language=Raden]
"../nativelibs/math" __sharedlib_ext to_string swap pop append loadnative
\end{lstlisting}

\section{Standard Library}

The standard library is stored in \code{libs/}. Functions are generally
loaded with \code{load} and called with \code{call}.

\subsection{core}

\code{libs/core.rdn} provides aliases and general helpers:

\begin{itemize}
\item Operator aliases: \code{not}, \code{add}, \code{sub}, \code{mul}, \code{div}.
\item Comparisons: \code{eq}, \code{neq}, \code{lt}, \code{gt}, \code{le}, \code{ge}.
\item Boolean and bitwise helpers: \code{and}, \code{or}, \code{xor}, \code{shl}, \code{shr}.
\item Type predicates: \code{is\_null}, \code{is\_int}, \code{is\_float}, \code{is\_bool}, \code{is\_str}, \code{is\_list}.
\item Numeric helpers: \code{inc}, \code{dec}, \code{neg}, \code{abs}, \code{even}, \code{odd}, \code{min}, \code{max}, \code{between}, \code{clamp}.
\item Functional helpers: \code{pipe}, \code{tap}, \code{select}, \code{when}, \code{unless}, \code{times}, \code{range}, \code{in}.
\end{itemize}

\subsection{lists}

\code{libs/lists.rdn} provides list helpers:
\code{each}, \code{foreach}, \code{map}, \code{collect}, \code{filter},
\code{reduce}, \code{copy}, \code{reverse}, \code{concat}, \code{first},
\code{head}, \code{last}, \code{tail}, \code{rest}, \code{count},
\code{isEmpty}, \code{includes}, \code{any}, \code{all}, \code{shift},
\code{popLast}, \code{unpack}, and \code{insert}.

\begin{lstlisting}[language=Raden]
(1 2 3) dup collect call print
(1 2 3) 2 > filter call print
(1 2 3) 0 add reduce call print
\end{lstlisting}

\subsection{strings}

\code{libs/strings.rdn} loads the native strings module and provides
\code{concat}, \code{concatBy}, \code{contains}, \code{hasPrefix},
\code{hasSuffix}, \code{trimSpace}, \code{split}, and \code{replaceAll}.

\subsection{io}

\code{libs/io.rdn} provides \code{readLine}, \code{readln},
\code{readAll}, \code{prompt}, \code{nl}, \code{println}, \code{say},
\code{printAll}, \code{printLines}, and \code{printJoin}.

\subsection{files}

\code{libs/files.rdn} provides \code{readLines}, \code{writeLines},
\code{appendLines}, \code{readText}, \code{writeText},
\code{appendText}, \code{slurp}, and \code{spit}.

\begin{lstlisting}[language=Raden]
"./libs/files.rdn" load
"./data.txt" readLines call lines let
lines each call print
\end{lstlisting}

\subsection{math}

\code{libs/math.rdn} provides \code{sqrt}, \code{powf}, and \code{pow}.
It also exposes constants such as \code{M\_E}, \code{M\_PI},
\code{M\_LN2}, \code{M\_SQRT2}, and related C math constants.

\subsection{typecheck}

\code{libs/typecheck.rdn} provides optional runtime signature checks for
signed \code{defun}. Load the library before defining signed functions.
The signed form is postfix: function name, parameter type list, return
type list, then \code{defun}.

\begin{lstlisting}[language=Raden]
"./libs/typecheck.rdn" load

add_ints (Int Int) (Int) defun
  +
end

2 3 add_ints check print
2 3 add_ints scall print
\end{lstlisting}

Use \code{()} when a function accepts no parameters or returns no values:

\begin{lstlisting}[language=Raden]
make_num () (Int) defun
  42
end

print_num (Int) () defun
  print
end
\end{lstlisting}

\begin{longtable}{ll}
\toprule
Name & Description \\
\midrule
\code{Int Str Bool Real List Null} & Concrete type constants for signatures \\
\code{Any} & Wildcard type; accepts any runtime value \\
\code{Func} & Function identifier type; matches values whose \code{type} is \code{"function"} \\
\code{check} & Validate arguments for a signed function without calling it \\
\code{scall} & Validate arguments, call through \code{pcall}, then validate returns \\
\bottomrule
\end{longtable}

\begin{lstlisting}[language=Raden]
identity_any (Any) (Any) defun
end

call_with_9 (Func) (Any) defun
  9 swap scall
end
\end{lstlisting}

\subsection{os, unix, path, process, net, json, strconv, flag}

\begin{longtable}{ll}
\toprule
Library & Main functions \\
\midrule
\code{os} & \code{pid}, \code{env}, \code{sleep}, \code{now} \\
\code{unix} & \code{pwd}, \code{exists}, \code{mkdir}, \code{rmPath}, \code{ls} \\
\code{path} & \code{join}, \code{basename}, \code{dirname}, \code{extname}, \code{absolute} \\
\code{process} & \code{status}, \code{output} \\
\code{net} & \code{host}, \code{encodeUrl}, \code{decodeUrl} \\
\code{json} & \code{parse}, \code{stringify} \\
\code{strconv} & \code{atoi}, \code{atof}, \code{parseBool}, \code{itoa}, \code{formatBool} \\
\code{flag} & \code{resetFlags}, \code{addFlag}, \code{execute}, \code{executeOn}, \code{positionals}, \code{programName}, \code{printHelp} \\
\bottomrule
\end{longtable}

\section{Native Modules}

Native modules are shared libraries written in C. They export:

\begin{lstlisting}[language=C]
bool rdn_module_init(RDNModule *module);
\end{lstlisting}

\subsection{Bundled Native Modules}

\begin{longtable}{ll}
\toprule
Module & Registered functions \\
\midrule
\code{files} & \code{readLines}, \code{readText}, \code{writeLines}, \code{appendLines}, \code{writeText}, \code{appendText} \\
\code{io} & \code{readLine}, \code{readAllInput} \\
\code{json} & \code{parseJson}, \code{stringifyJson} \\
\code{math} & \code{powerOf}, \code{sqrtOf} \\
\code{net} & \code{hostName}, \code{urlEncode}, \code{urlDecode} \\
\code{path} & \code{joinPath}, \code{baseName}, \code{dirName}, \code{extName}, \code{isAbsolutePath} \\
\code{process} & \code{execStatus}, \code{execOutput} \\
\code{strconv} & \code{atoi}, \code{atof}, \code{parseBool}, \code{itoa}, \code{formatBool} \\
\code{strings} & \code{strContains}, \code{strHasPrefix}, \code{strHasSuffix}, \code{strTrimSpace}, \code{strSplit}, \code{strReplaceAll} \\
\code{syscall} & \code{pid}, \code{getEnv}, \code{sleepMs}, \code{epochTime} \\
\code{unix} & \code{cwd}, \code{pathExists}, \code{makeDir}, \code{removePath}, \code{listDir} \\
\bottomrule
\end{longtable}

\subsection{Native Stack API}

\begin{longtable}{ll}
\toprule
Function & Description \\
\midrule
\code{api->stack\_size(api)} & Number of items on the stack \\
\code{api->type(api, index)} & Type of value at index; \code{-1} is top \\
\code{api->to\_integer(api, index, \&out)} & Read integer \\
\code{api->to\_number(api, index, \&out)} & Read double \\
\code{api->to\_boolean(api, index, \&out)} & Read boolean \\
\code{api->to\_string(api, index)} & Read string as \code{const char *} \\
\code{api->pop(api, count)} & Pop values \\
\code{api->push\_integer(api, value)} & Push integer \\
\code{api->push\_number(api, value)} & Push double \\
\code{api->push\_boolean(api, value)} & Push boolean \\
\code{api->push\_string(api, value)} & Push copied string \\
\code{api->raise\_error(api, message)} & Raise an error and return false \\
\bottomrule
\end{longtable}

\subsection{Native List API}

\begin{longtable}{ll}
\toprule
Function & Description \\
\midrule
\code{api->push\_list(api)} & Push an empty list \\
\code{api->list\_len(api, index, \&len)} & Read list length \\
\code{api->list\_append(api, list\_idx, value\_idx)} & Append cloned value \\
\code{api->list\_index(api, list\_idx, item\_idx)} & Push cloned item \\
\code{api->list\_remove(api, list\_idx, item\_idx)} & Remove item \\
\bottomrule
\end{longtable}

\subsection{Native Module Example}

\begin{lstlisting}[language=C]
#include "../include/rdn_native.h"

static bool native_add(RDNApi *api) {
  long left = 0, right = 0;

  if (api->stack_size(api) < 2)
    return api->raise_error(api, "add requires 2 operands");

  if (!api->to_integer(api, -2, &left) ||
      !api->to_integer(api, -1, &right))
    return api->raise_error(api, "add requires integers");

  api->pop(api, 2);
  return api->push_integer(api, left + right);
}

bool rdn_module_init(RDNModule *module) {
  return module->register_function(module, "add", native_add);
}
\end{lstlisting}

\begin{lstlisting}[language=bash]
gcc -Wall -Wextra -Werror -ggdb -std=c11 -fPIC -shared \
  mymodule.c -I. -o mymodule.so
\end{lstlisting}

\begin{lstlisting}[language=Raden]
"./mymodule.so" loadnative
3 4 add call print
\end{lstlisting}

\section{Embedding}

Raden exposes a minimal embedding API through \code{include/rdn.h}.

\begin{lstlisting}[language=C]
#include "./include/rdn.h"

int main(void) {
  char *argv[] = {
    "host",
    "./script.rdn",
  };
  return rdn_main(2, argv);
}
\end{lstlisting}

The first argument is the program name and the second is the script path.
If \code{argc < 2}, \code{rdn\_main} starts the REPL. Source-relative
\code{load} and \code{loadnative} path resolution work in embedded usage.

\section{Editor Support}

The repository includes Neovim runtime files:

\begin{lstlisting}[language=bash]
cp -r ./nvim/* ~/.config/nvim/
\end{lstlisting}

This provides file detection for \code{.rdn}, syntax highlighting,
ftplugin defaults, and basic indent rules. A starter Treesitter grammar
is available in \code{treesitter/raden/}.

\section{Test Suite}

Raden includes regression tests in \code{test/}. Each test is a \code{.rdn}
file with expected output in \code{test/expected/}. The runner compares
stdout and stderr and verifies the exit code.

\begin{lstlisting}[language=bash]
bash test/run.sh
\end{lstlisting}

Useful tests include:
\code{test/bugs.rdn}, \code{test/lists.rdn}, \code{test/loop.rdn},
\code{test/functions.rdn}, \code{test/strings\_and\_type.rdn}, and
\code{test/rule\_110.rdn}.

\section{Implementation Notes}

When editing the interpreter:

\begin{enumerate}
\item Remember that the language is postfix; many apparent bugs are
operand-order mistakes.
\item List literals execute code and collect resulting stack values.
\item \code{type} recognizes function identifiers.
\item Identifier handling is not completely uniform because some values
are cloned while list and string names may be preserved for mutation.
\item Validate changes with \code{make} and \code{bash test/run.sh}.
\end{enumerate}

\end{document}
